\begin{afterword}
\summaryhead

The post opens with two epigraphs from E. E. Smith's Lensman novels about reading a whole universe from one pebble. The author doubts that a pebble fixes the facts of human history, but holds that it carries our laws of physics and marks of the planet it came from, and that a geologist might catch someone lying about where a pebble was found.

Everything is physically connected to its past; a penny even pulls on the Moon. Only some connections can be noticed, and these form what the author calls the ``Great Web of Causality.'' A lie is a fact out of place in this web. Defending it can force further lies about connected facts, as in a short dialogue about a business trip, and people often fail to foresee the facts that will expose them: genetics, later forensic evidence, perhaps a future lie detector or Paul Ekman. Not all lies are caught, the author grants, but lies are riskier than liars think, and honesty or silence carries less unforeseen risk.

\responsehead

The post is careful about what it claims. It says twice that not all lies are uncovered, and once that not all lies spin out of control. It rejects the strong version of its own epigraphs. It ends by declining to call honesty the best policy, and claims only that honesty or silence carries less unforeseen risk than lying. Its geology is sound in outline: minerals do record the pressures they formed under, and composition is how some meteorites were traced to Mars.

The difficulty is that the first half of the post does not support the second. The physics shows that everything is connected, and the paragraph on a penny's pull on the Moon shows that most connections leave no usable trace. Whether a lie is exposed therefore depends on which traces are noticeable and on whether anyone checks. That is a question about lies and liars, and for it the post offers a made-up dialogue and a list of general examples. The claim that the web is ``very commonly underestimated'' is a claim about what liars expect, and no evidence about liars is given.

The one specific claim about detecting lies goes beyond the evidence of its time. Paul Ekman, the post says, ``could probably read off a sizeable fraction of the world's lies.'' The research up to 2008 found that people detect lies about 54\% of the time, and was divided on whether some experts do much better. The post cites none of it.

Two smaller points. The penny figure fits a one-gram coin; a real cent gives about two and a half times as much, which does not change the point. And the claim that a pebble implies ``all the Everett branches'' assumes the many-worlds interpretation, which the post takes as given.

In short: a well-hedged case that lies are riskier than they look, whose physics shows connection but not detection, and whose one specific claim about detecting lies, on Paul Ekman, is given no support.
\end{afterword}
